\section{Modelo continuo}
\label{sec:modelo-continuo}

\subsection{Dominio y supuestos físicos}
La fundación crece durante la colocación. Si $\Omega_j$ representa la capa $j$ y $t_j$ su instante de activación, el dominio de conducción y la edad local se expresan como
\begin{equation}
  \Omega_c(t)=\bigcup_{j:\,t\geq t_j}\Omega_j,
  \qquad a_j(t)=t-t_j.
  \label{eq:dominio-edad}
\end{equation}
Solo el hormigón colocado almacena y conduce calor. Su densidad $\rho$, su calor específico $c_p$ y su conductividad $k_c$ se suponen constantes e iguales en todas las capas. Esta elección permite aislar el efecto del cronograma con las propiedades disponibles. El terreno, el moldaje y el agua se representan por sus intercambios con el hormigón.

\subsection{Balance local de energía y conducción}
En un volumen fijo $V$ de hormigón activo, el cambio de energía sensible es igual al calor que atraviesa el contorno más el generado por hidratación y menos el retirado por los circuitos \citep{inlTransferencia}. Con normal saliente $\vec n$, flujo conductivo $\vec q$ en $\mathrm{W\,m^{-2}}$ y temperatura de referencia constante $T_{\mathrm{ref}}$, el balance y la ley de Fourier son
\begin{equation}
  \frac{d}{dt}\int_V \rho c_p(T-T_{\mathrm{ref}})\,dV
  =-\int_{\partial V}\vec{q}\cdot\vec{n}\,dA
   +\int_V(\dot q_h-\dot q_r)\,dV.
  \label{eq:balance-integral}
\end{equation}
\begin{equation}
  \vec{q}=-k_c\vec{\nabla} T.
  \label{eq:fourier}
\end{equation}
Al aplicar el teorema de la divergencia al flujo y considerar un volumen arbitrario se obtiene la ecuación que se resuelve dentro del hormigón:
\begin{equation}
  \rho c_p\frac{\partial T}{\partial t}
  =\vec{\nabla}\cdot(k_c\vec{\nabla} T)+\dot q_h-\dot q_r,
  \qquad \vec{x}\in\Omega_c(t).
  \label{eq:calor-modelo}
\end{equation}
Los términos $\dot q_h$ y $\dot q_r$ expresan potencia por unidad de volumen, en $\mathrm{W\,m^{-3}}$. El primero calienta el material y el segundo retira energía cuando el sistema de refrigeración lo permite.

\subsection{Colocación, condiciones iniciales e hidratación}
Para una tasa constante $q$ y sin esperas, cada capa de volumen $V_j$ tarda $V_j/q$ en colocarse. El modelo activa todo su volumen al comienzo de ese intervalo. El hormigón ya colocado conserva su temperatura y el nuevo entra con la condición
\begin{equation}
  T(\vec{x},t_j^+)=T_{\mathrm{col}}=12\,^{\circ}\mathrm{C},
  \qquad \vec{x}\in\Omega_j.
  \label{eq:inicial-colocacion}
\end{equation}
Todas las capas de un escenario utilizan la misma curva $m$, desplazada según su edad local. La fuente efectiva se obtiene del calor acumulado en la sección anterior:
\begin{equation}
  \dot q_h(\vec{x},t)=\frac{dH_m}{da}\bigl(a_j(t)\bigr),
  \qquad \vec{x}\in\Omega_j,\quad t\geq t_j.
  \label{eq:fuente-edad}
\end{equation}
La velocidad de hidratación depende solo de la edad. En un ensayo adiabático bastaría multiplicar la tasa de aumento de temperatura por $\rho c_p$. Las probetas docentes pierden calor por su superficie. El balance empleado para interpretar la inversión incluye también esa pérdida:
\begin{equation}
  \dot q_m(t)=\frac{1}{V_p}\frac{d}{dt}
       \int_{V_p}\rho c_p(T_p-T_{\mathrm{ref}})\,dV
       +\frac{1}{V_p}\int_{\partial V_p}\beta_p(T_{p,s}-T_a)\,dA.
  \label{eq:balance-probeta}
\end{equation}
Aquí $V_p$ es el volumen de la probeta, $T_p$ su temperatura interior, $T_{p,s}$ su temperatura superficial y $T_a$ la del aire. La conductancia $\beta_p=2{,}5\,\mathrm{W\,m^{-2}\,K^{-1}}$ corresponde al ajuste docente. El modelo de la fundación recibe el resultado de ese ajuste.

\subsection{Refrigeración equivalente y control}
La refrigeración actúa únicamente en la región de cada circuito abierto. La extracción aumenta con el exceso de temperatura del hormigón sobre el agua, que se supone uniforme a $T_w=18\,^\circ\mathrm C$:
\begin{equation}\begin{aligned}
  \dot q_r^0(\vec{x},t)&=U_v\sum_{b=1}^{9}
       \chi_b(\vec{x})\,s_b(t)\max(T-T_w,0),\\
  P_r^0(t)&=\int_{\Omega_c(t)}\dot q_r^0\,dV,
  \qquad \dot q_r=\lambda(t)\dot q_r^0,
  \end{aligned}\label{eq:extraccion-equivalente}\end{equation}
La función $\chi_b$ identifica la región del ramal $b$ y $s_b$ vale uno si está abierto. El coeficiente volumétrico $U_v$ se fija en $5\,\mathrm{W\,m^{-3}\,K^{-1}}$. Si la potencia solicitada $P_r^0$ supera $15\,\mathrm{kW}$, se reducen todos los aportes en la misma proporción mediante $\lambda=P_{\max}/P_r^0$; en los demás casos $\lambda=1$.

Cada circuito comienza cerrado. Una vez cubierto, abre cuando el punto más frío de sus controles alcanza $19\,^\circ\mathrm C$ y cierra cuando baja a $18\,^\circ\mathrm C$ o menos. La ley representa la capacidad de extracción adoptada. Determinar si el montaje la alcanza exige comprobar la transferencia real, la hidráulica y el calentamiento del agua a lo largo del ramal.

\subsection{Superficies expuestas y protecciones}
La madera y las mantas se representan como resistencias en serie. El calor entrante $\phi=-\vec q\cdot\vec n$ es positivo hacia el hormigón. Entre la cara del hormigón a $T_b$ y el exterior de la protección a $T_s$ se cumple
\begin{equation}
  R_p(a)=\sum_{\ell\,\mathrm{instalada}}\frac{e_\ell}{k_\ell},
  \qquad T_s-T_b=R_p(a)\,\phi,
  \label{eq:resistencia-proteccion}
\end{equation}
Los espesores $e_\ell$ y las conductividades $k_\ell$ corresponden a las capas instaladas en ese instante. Se supone despreciable su almacenamiento de calor y la resistencia de sus contactos. Si solo hubiera convección, se obtendría el coeficiente equivalente
\begin{equation}
  -k_c\vec{\nabla} T\cdot\vec{n}=h_{\mathrm{eq}}(T_b-T_a),
  \qquad h_{\mathrm{eq}}^{-1}=h_c^{-1}+R_p.
  \label{eq:robin-equivalente}
\end{equation}
En el cálculo se resuelve además el intercambio radiativo en la cara exterior. Al retirar una protección se elimina su resistencia. Cuando se retira la última, $T_s=T_b$.

\subsection{Radiación solar y de onda larga}
La cara exterior comparte calor con el aire, absorbe radiación solar y emite radiación de onda larga hacia el cielo y el entorno \citep{energyplusMcAdams}. Su balance por unidad de superficie es
\begin{equation}
  \phi=\mathcal F(T_s)
  =h_c(T_a-T_s)+\alpha_{\mathrm{SW}}I_f
   +\varepsilon\sigma\left[
      F_{\mathrm{sky}}T_{\mathrm{sky},K}^{4}
      +(1-F_{\mathrm{sky}})T_{a,K}^{4}-T_{s,K}^{4}\right].
  \label{eq:balance-exterior}
\end{equation}
$T_a$ representa la temperatura del aire, $I_f$ la irradiancia que recibe la cara y $T_{\mathrm{sky},K}$ la temperatura radiativa del cielo. El viento determina $h_c$, mientras que $\alpha_{\mathrm{SW}}$, $\varepsilon$ y $\sigma$ son la absortancia, la emisividad y la constante de Stefan--Boltzmann. Las temperaturas elevadas a la cuarta potencia se expresan en kelvin. El factor de vista $F_{\mathrm{sky}}$ vale $1$ arriba, $0{,}5$ en caras verticales y $0$ abajo. La temperatura efectiva del cielo se estima con el punto de rocío y la nubosidad, mediante las relaciones de Clark--Allen y Walton \citep{energyplusCielo}.

El eje $+x$ apunta al este y el eje $+y$ al norte. Las caras superiores reciben la irradiancia global horizontal. En las verticales se suman la radiación directa proyectada sobre la normal, la mitad de la difusa y la mitad de la global reflejada por el terreno con albedo $0{,}20$. Las inferiores expuestas reciben solo la componente reflejada. La posición del sol se calcula a mitad de cada hora histórica \citep{noaaSolar}. Desde la edad local de $504\,\mathrm h$, la sombra multiplica la irradiancia por $0{,}10$ y conserva la convección y la onda larga.

\subsection{Suelo, interfaces y cambios de borde}
La base conduce calor a través de una resistencia equivalente $R_g=L_g/k_g$ hacia un reservorio que permanece a $T_g=15\,^\circ\mathrm C$ desde el inicio. El modelo omite la evolución térmica del suelo. Entre capas activas se supone contacto perfecto, con una misma normal $\vec n_j$ para comparar los flujos a ambos lados:
\begin{equation}
  -k_c\vec{\nabla} T\cdot\vec{n}=\frac{T_b-T_g}{R_g}.
  \label{eq:borde-suelo}
\end{equation}
\begin{equation}
  T^-=T^+,\qquad
  -k^-\vec{\nabla} T^-\cdot\vec{n}_j
  =-k^+\vec{\nabla} T^+\cdot\vec{n}_j.
  \label{eq:contacto-capas}
\end{equation}
Cuando una capa nueva cubre una cara, el intercambio con el aire cesa y comienza la conducción hacia la capa vecina. La Tabla~\ref{tab:bordes-etapas} reúne las condiciones aplicadas ($R_m=e_m/k_m$ para la madera y $R_a=e_a/k_a$ para cada manta), incluidas las etapas durante las cuales la madera o las mantas siguen protegiendo cada superficie.
\begin{table}[htbp]
  \centering
  \caption{Condiciones de borde e interfaces por superficie y etapa}
  \label{tab:bordes-etapas}
  \small
  \setlength{\tabcolsep}{3pt}
  \begin{tabular}{L{2.5cm}L{2.0cm}L{3.7cm}L{3.0cm}L{4.5cm}}
    \toprule
    Superficie & Etapa & Ecuación & Parámetros & Fuente y fundamento \\
    \midrule
    Lateral al aire & $0\leq a<72\,\mathrm h$ & $\phi=\mathcal F(T_s)$ & $R_p=R_m+2R_a$ & AraucoPly, Aislanglass y plan: madera con dos mantas \\
    Lateral al aire & $72\leq a<240\,\mathrm h$ & $\phi=\mathcal F(T_s)$ & $R_p=2R_a$ & Plan: retiro de madera y conservación de mantas \\
    Superior al aire & $0\leq a<240\,\mathrm h$ & $\phi=\mathcal F(T_s)$ & $R_p=2R_a$ & Aislanglass y plan: dos mantas, sin madera \\
    Lateral o superior al aire & $240\leq a<504\,\mathrm h$ & $\phi=\mathcal F(T_s)$ & $R_p=R_a$ & Plan: retiro de la primera manta \\
    Lateral o superior al aire & $a\geq504\,\mathrm h$ & $\phi=\mathcal F(T_b)$ con $I_f$ atenuada & $R_p=0$, $\tau=0{,}10$ & Plan: retiro de última manta y sombra idealizada \\
    Base sobre suelo & Desde su activación & $\phi=(T_g-T_b)/R_g$ & $T_g=15\,^{\circ}\mathrm C$, $R_g=L_g/k_g$ & Supuesto de reservorio y conducción equivalente \\
    Interfaz interna & Ambas capas activas & $T^-=T^+$ y continuidad de flujo, Ec.~\eqref{eq:contacto-capas} & Resistencia de junta nula & Contacto perfecto supuesto, sin borde de aire \\
    \bottomrule
  \end{tabular}
  \captionsetup{skip=2pt}
  \caption*{Fuente: configuración final del estudio, referencias de materiales y ecuaciones del modelo implementado.}
\end{table}
